\uuid{If44}
\exo7id{3665}
\titre{exo7 3665}
\auteur{quercia}
\organisation{exo7}
\datecreate{2010-03-11}
\isIndication{false}
\isCorrection{true}
\chapitre{Espace euclidien, espace normé}
\sousChapitre{Produit scalaire, norme}
\cpge{mpsi-mp2i.espaces-prehilbertiens.b, pcsi.espaces-prehilbertiens.b, ptsi.geometrie-espace.b, mpsi-mp2i.espaces-prehilbertiens.c, pcsi.espaces-prehilbertiens.c, mpsi-mp2i.espaces-prehilbertiens.e, pcsi.espaces-prehilbertiens.e, psi.espaces-prehilbertiens.b, pt.structure-prehilbertienne.c, psi.espaces-prehilbertiens.d, pt.structure-prehilbertienne.e, bcpst2.produit-scalaire.a, bcpst2.produit-scalaire.b}
\module{Algèbre}
\niveau{L2}
\difficulte{}

\contenu{
\texte{
On munit $E = \R_n[X]$ du produit scalaire :
Pour $P = \sum_i a_iX^i$ et $Q = \sum_i b_iX^i$,
$(P\mid Q) = \sum_i a_ib_i$.

Soit $H = \{ P \in E \text{ tq } P(1) = 0 \}$.
}
\begin{enumerate}
    \item \question{Trouver une base orthonormale de $H$.}
    \item \question{Calculer $d(X,H)$.}
\reponse{
$\frac1{\sqrt{n+1}}$.
}
\end{enumerate}
}
